% SPDX-License-Identifier: GPL-3.0-or-later
% Copyright (C) 2026 Alessandro Greco
\chapter{Академические исследования Tor}\label{cap:29-ricerca}
\aggiornatoal{23 сентября 2026 г.}{работы, цитируемые в этой главе; исследовательский портал Tor Project; консенсус от 23 сентября 2026 г.}

\begin{obiettivi}
К концу этой главы вы сможете:
\begin{itemize}[nosep]
  \item находить основные площадки публикации исследований Tor и ресурсы Tor Project;
  \item количественно измерять анонимность — от энтропии до времени первого компрометационного события;
  \item выбирать подходящий экспериментальный метод: аналитическое моделирование, симуляцию, эмуляцию или измерение;
  \item симулировать тысячи клиентов на протяжении месяцев, используя реальный консенсус (Лабораторная работа~19).
\end{itemize}
\end{obiettivi}

Tor с самого начала развивался под влиянием исследований. Onion routing появился в исследовательской лаборатории (\cref{cap:03-nascita-tor}); архитектура Tor была представлена на академической конференции~\cite{dingledine2004tor}; многие механизмы, рассматриваемые в книге, от сторожевых узлов до управления перегрузкой, возникли как ответ на опубликованные исследования. Эта глава предназначена для читателей, которые хотят заниматься исследованиями сами: где публиковаться, как измерять анонимность, какие экспериментальные инструменты использовать и какие принципы безопасности соблюдать.

% ---------------------------------------------------------------------------
\section{Где публикуются исследования Tor}\label{sec:ric-dove}

Исследования сетей анонимности в основном публикуются на площадках двух типов:
\begin{description}
  \item[Privacy Enhancing Technologies Symposium (PETS)] и его журнал \emph{Proceedings on Privacy Enhancing Technologies} (PoPETs) — специализированная площадка; ключевая работа по измерению анонимности, рассматриваемая далее, была представлена на семинаре PET в 2002 году, до появления журнала PoPETs~\cite{diaz2002measuring}.
  \item[Крупные конференции по безопасности] USENIX Security, ACM CCS, IEEE Symposium on Security and Privacy и NDSS. Многие статьи из \cref{cap:26-attacchi} были опубликованы именно там.
\end{description}
\emph{Anonymity Bibliography} проекта Free Haven более двадцати лет собирает публикации в этой области и является естественной отправной точкой для обзора литературы.

Tor Project также поддерживает исследовательский портал с данными, инструментами, техническими отчётами, предлагаемыми темами исследований и информацией о Research Safety Board. Рекомендации подчёркивают, что протокол и сеть постоянно меняются: измерения, выполненные без понимания внутренней работы, могут приводить к неверным выводам. Поэтому исследователям рекомендуется связываться с проектом заранее, чтобы узнать о параллельных экспериментах и планируемых изменениях.

% ---------------------------------------------------------------------------
\section{Измерение анонимности}\label{sec:ric-metriche}

В \cref{cap:01-anonimato} анонимность определялась через неразличимость внутри множества. Однако с точки зрения атакующего не каждый участник обязательно одинаково вероятен. После наблюдения трафика противник может назначить некоторым пользователям гораздо более высокие вероятности, чем другим. Одного размера множества недостаточно.

\paragraph{Степень анонимности.} В 2002 году Díaz и соавторы предложили измерять анонимность через энтропию распределения вероятностей, которое атакующий назначает возможным отправителям~\cite{diaz2002measuring}. Если $p_i$ — вероятность, приписанная пользователю $i$ как отправителю, энтропия равна
\[
H(X) = -\sum_{i=1}^{N} p_i \log_2 p_i,
\]
и достигает максимума, $H_M = \log_2 N$, когда все $N$ пользователей равновероятны. \textbf{Степень анонимности} определяется как
\[
d = \frac{H(X)}{H_M}, \qquad 0 \le d \le 1:
\]
при $N\ge 2$ она равна 1 для равномерного апостериорного распределения (это не обязательно означает отсутствие новых сведений), а 0 — когда вероятность одного отправителя равна 1. При $N=1$ отношение не определено, поскольку $H_M=0$.

Интуитивная интерпретация энтропии — $2^{H(X)}$: число равновероятных пользователей, создающее ту же неопределённость. Если среди 1\,000 пользователей атакующий сузил подозреваемых до двух с вероятностью 0,5 для каждого, неопределённость равна множеству из двух: $H = 1$ бит.

\paragraph{Ограничения разового измерения.} Энтропия описывает отдельное наблюдение, но экспозиция Tor накапливается во времени, как мы видели в \cref{sec:att-correlazione}: каждая цепь создаёт новую возможность для атакующего. Поэтому Johnson и коллеги предложили \emph{продольные} метрики, включая вероятность хотя бы одной деанонимизации за заданный период и медианное время до первого компрометационного события, вычисляемые путём симуляции поведения пользователей на реальных сетевых данных~\cite{johnson2013users}. В лабораторной работе мы используем сходный подход.

% ---------------------------------------------------------------------------
\section{Четыре метода}\label{sec:ric-metodi}

\begin{description}
  \item[Аналитическое моделирование] Математическая модель, например степень анонимности или модель Sybil-атак из \cref{sec:att-sybil}. Метод быстрый и общий, но зависит от предположений, которые реальные системы в той или иной степени нарушают.
  \item[Поведенческая симуляция] Воспроизводится упрощённый компонент системы, например выбор маршрута, на основе реальных данных. TorPS, симулятор маршрутов Джонсона и соавторов~\cite{johnson2013users}, и наша Лабораторная работа~19 используют этот подход. Он быстрый, но не учитывает ничего, что не включено в модель.
  \item[Сетевая эмуляция и симуляция] Реальное ПО запускается либо в приватной сети, например лаборатории Chutney, используемой в книге, либо внутри дискретно-событийного симулятора вроде Shadow~\cite{jansen2012shadow}, способного запускать тысячи ретрансляторов и клиентов на одном компьютере с виртуальным временем (\cref{sec:rp-shadow}). Особенно полезно для исследований производительности: предлагаемые изменения протокола можно проверить в реалистичной сети до развёртывания.
  \item[Измерения в публичной сети] Только прямое измерение показывает, что происходит в реальной сети. Оно же создаёт наибольший риск для реальных пользователей и требует особенно осторожного дизайна эксперимента (\cref{sec:ric-sicurezza}).
\end{description}
Сильные исследования часто сочетают несколько методов: математический анализ формулирует гипотезу, симуляция проверяет её в масштабе, а осторожно спроектированное измерение оценивает реалистичность.

\begin{approfondimento}[Измерение без наблюдения отдельных пользователей]
PrivCount, разработанный Jansen и Johnson, показывает, как измерения могут сочетаться с приватностью~\cite{jansen2016privcount}. Участвующие ретрансляторы подсчитывают нужные события, но счётчики маскируются случайными значениями, разделяемыми с другими участниками. Добавляется шум с применением принципов \emph{дифференциальной приватности}. Видимым становится только агрегированный, намеренно зашумлённый результат, что ограничивает выводы об отдельных пользователях. В 2016 году авторы с помощью PrivCount оценили примерно 710\,000 одновременно подключённых пользователей, из которых 550\,000 были активны в данный момент, и установили, что на трафик к веб-портам приходился 91\% байтов данных, наблюдавшихся на участвующих выходных узлах. Это оценка исследования 2016 года, а не характеристика современной сети. Это развивает те же принципы приватности, что и округление статистики ретрансляторов (\cref{sec:metrics-utenti}).
\end{approfondimento}

% ---------------------------------------------------------------------------
% Колонтитул короче: полное название касается номера страницы. Для правого колонтитула
% LaTeX берёт первую метку на странице, поэтому метку самого \section отключаем.
{\renewcommand{\sectionmark}[1]{}%
\section{Лабораторная работа: время до первого компрометационного события}\label{sec:lab-ricerca}}
\sectionmark{Лабораторная работа: время до компрометации}

\begin{laboratorio}[19]{Месяцы цепей в симуляции}
\textbf{Цель.} Используя реальный консенсус, оценить, сколько клиентов атакующий, контролирующий заданную долю пропускной способности, может деанонимизировать со временем при постоянных guard и без них.\\
\textbf{Предварительные условия.} Лабораторные работы~7 и~15; работающая сеть не требуется.\\
\textbf{Оценочное время.} 1--2 часа.
\end{laboratorio}

Модуль \file{labs/ru/код/torlab/simulation.py} повторно использует веса из Лабораторной работы~7. Атакующий добавляет сторожевой и выходной ретрансляторы, каждый из которых представляет долю $q$ взвешенной пропускной способности в своей позиции цепи. Мы моделируем 2\,000 клиентов в течение 180 дней, каждый строит по 20 цепей в день.

В модели постоянный сторожевой узел меняется по приближённому расписанию для \emph{подтверждённого} сторожевого узла. Согласно guard-spec и коду \file{entrynodes.c}, подтверждённый сторожевой узел удаляется из выборки через \texttt{GUARD\_CONFIRMED\_MIN\_LIFETIME} (60 дней) после подтверждения, а \texttt{CONFIRMED\_ON\_DATE} случайно сдвигается назад максимум на \texttt{GUARD\_LIFETIME}/10 (12 дней). Поэтому в модели период смены узла составляет 48--60 дней. Параметр \texttt{GUARD\_LIFETIME}, равный 120 дням, относится к сроку хранения неподтверждённых узлов в выборке. Настоящий клиент учитывает также списки узлов, доступность и замены: модель не воспроизводит полный алгоритм. Ранняя версия этой лабораторной работы ошибочно моделировала ротацию через 120 дней и занижала результат для выбранной упрощённой модели. Для сравнения мы также моделируем клиентов, выбирающих новый guard для каждой цепи. Клиент считается скомпрометированным в первый день, когда одна из его цепей содержит и вредоносный guard, и вредоносный exit.

Модель намеренно проста: она игнорирует промежуточные ретрансляторы, семейства, сети /16, резервные сторожевые узлы и отказы. Ради скорости она не использует \texttt{random.choices}, который пересчитывал бы тысячи весов при каждом выборе. Вместо этого используется бинарный поиск по накопленным весам, позволяющий выполнить миллионы выборок без повторного расчёта всех весов; реальное время зависит от компьютера.

Для воспроизведения сначала загрузите исторический открытый консенсус (\texttt{make -C labs/ru data}): файл не хранится в Git. Симуляция детерминирована: при той же модели, консенсусе, доле и зерне результат всегда один и тот же:
\begin{lstlisting}[language=bash,numbers=none]
$ PYTHONPATH=labs/ru/код python3 -m torlab.simulation labs/ru/данные/collector/consensus-microdesc-2026-09-23-15 0.05
противник: 5% пропускной способности сторожевых узлов и 5% выходных узлов; 2000 клиентов, 180 дней, 20 цепей/день
постоянный сторожевой узел    1 дн.:  2.6%    7 дн.:  4.4%   30 дн.:  4.4%   90 дн.:  9.4%  180 дн.: 18.4%  медиана: None
новый сторожевой узел на цепь 1 дн.:  4.3%    7 дн.: 27.9%   30 дн.: 76.5%   90 дн.: 98.9%  180 дн.:100.0%  медиана: 15
\end{lstlisting}

\begin{figure}[htbp]
\centering
% SPDX-License-Identifier: GPL-3.0-or-later
% Copyright (C) 2026 Alessandro Greco
%% Русское альтернативное описание EPUB.
\altfig{Две накопительные кривые учебной симуляции лабораторной работы 19: 2000 клиентов в течение 180 дней, по 20 цепей в день. Атакующий контролирует 5\% взвешенной полосы сторожевых и 5\% выходных ретрансляторов. Компрометация в модели означает хотя бы одну цепь с обоими вредоносными концами: к 180-му дню 100\% при новом guard для каждой цепи и 18,4\% при постоянном guard. Ротации распределены по клиентам с интервалами 48--60 дней, а не синхронизированы. Это результаты симуляции, а не наблюдения реальной деанонимизации.}
% Ротации для каждого клиента независимы (48--60 дней в модели), без общих дат.
\begin{tikzpicture}[font=\sffamily\scriptsize,x=0.04cm,y=0.04cm]
  \draw[->] (0,0) -- (190,0) node[below left] {дни};
  \draw[->] (0,0) -- (0,108) node[above] {\% скомпрометировано (модель)};
  \foreach \x in {0,30,60,90,120,150,180} \draw (\x,0) -- (\x,-2) node[below] {\x};
  \foreach \y in {0,25,50,75,100} \draw (0,\y) -- (-2,\y) node[left] {\y};
  \draw[thick,avviso] plot coordinates {(0,0.0) (5,21.1) (10,38.5) (15,51.1) (20,61.5) (25,69.7) (30,76.5) (35,81.7) (40,86.0) (45,89.2) (50,91.7) (55,93.3) (60,94.9) (65,96.2) (70,97.2) (75,97.9) (80,98.2) (85,98.7) (90,98.9) (95,99.1) (100,99.5) (110,99.7) (120,99.9) (165,100.0) (180,100.0)};
  \draw[thick,torviola] plot coordinates {(0,0.0) (1,2.6) (3,4.3) (5,4.4) (48,4.4) (52,5.8) (56,7.1) (60,8.7) (64,9.4) (100,9.5) (108,10.8) (114,12.4) (120,13.2) (130,13.4) (160,14.4) (165,15.6) (170,17.0) (175,18.2) (180,18.4)};
  \node[avviso,anchor=west] at (40,80) {новый guard для каждой цепи};
  \node[torviola,anchor=west] at (66,26) {постоянный guard};
% Ротации для каждого клиента независимы (48--60 дней в модели), без общих дат.
  \node[torgrigio,anchor=west,align=left] at (66,43) {ротации guard\\не синхронны};
\end{tikzpicture}

\caption{Лабораторная работа~19, упрощённая симуляция: 2000 клиентов, по 20 цепей в день в течение 180 дней, консенсус от 23 сентября 2026 года; атакующий контролирует 5\% взвешенной полосы сторожевых и 5\% выходных узлов. Накопительная доля отражает клиентов, чья хотя бы одна цепь в модели содержала оба вредоносных конца. Ротации постоянных guard распределены по клиентам (модельные интервалы 48--60 дней), а не происходят в три общие даты. Это не наблюдения реальной деанонимизации в публичной сети.}
\label{fig:compromissione}
\end{figure}

\Cref{fig:compromissione} показывает обе кривые. Без постоянных guard половина клиентов оказывается скомпрометирована за пятнадцать дней, а почти все — за три месяца, потому что каждая новая цепь создаёт ещё одну независимую возможность. При постоянных guard экспозиция клиента в основном определяется моментом выбора guard. Около 4\% клиентов, выбравших вредоносный guard, компрометируются рано; остальные не удовлетворяют этому конкретному условию компрометации, пока сохраняют честный сторожевой узел. По мере того как разные клиенты меняют guard в разные моменты времени (примерно каждые 48--60 дней для каждого клиента в модели), возникают новые возможности и доводит итог примерно до 18\% за шесть месяцев. При доле вредоносной пропускной способности 1\% соответствующий шестимесячный результат падает примерно до 4\% против 31\% без постоянных guard.

Эти результаты, основанные на эталонном снимке сентября 2026 года, иллюстрируют назначение сторожевых узлов. В этой упрощённой модели их главное преимущество проявляется в распределении риска между клиентами, а не в гарантированном снижении вероятности на каждой отдельной цепи. Для большинства пользователей смоделированный риск сначала близок к нулю, но для меньшей группы он высок. Поэтому время жизни guard — компромисс. Более медленная ротация дольше защищает пользователей с честным guard, но также продлевает экспозицию тех, кто выбрал вредоносный. Терпеливый атакующий, остающийся в сети месяцами, получает новые возможности при каждой ротации.

% ---------------------------------------------------------------------------
\section{Безопасное проведение исследований}\label{sec:ric-sicurezza}

В \cref{cap:25-legale} были введены рекомендации Tor Research Safety Board~\cite{trsb}. При проектировании эксперимента сначала спросите: \emph{действительно ли исследовательский вопрос требует публичной сети?} Многие вопросы производительности и атак можно изучить с помощью Shadow или приватной сети, получая воспроизводимые эксперименты и контроль переменных. Если измерение публичной сети действительно необходимо, следуйте уже обсуждавшимся принципам: воздействуйте только на собственный трафик, собирайте минимум данных, агрегируйте до публикации и заранее консультируйтесь с Board и Tor Project.

Недавняя история даёт противоположные примеры. В \cref{cap:27-casi-rete} описан эксперимент или атака, затронувшая реальных пользователей через сигнал, доступный получающим ретрансляторам, но не непосредственно стороннему наблюдателю за зашифрованными TLS-соединениями. PrivCount, напротив, показывает, что подробные измерения сети возможны без раскрытия отдельных пользователей. Разница заключается не только в инструментах, но и в дизайне эксперимента.

% ---------------------------------------------------------------------------
\section*{Итоги}
\begin{itemize}
  \item Исследования Tor публикуются на PETS и крупных конференциях по безопасности. Tor Project предоставляет данные, инструменты и Research Safety Board.
  \item Энтропия измеряет анонимность для отдельного наблюдения; степень анонимности нормализует её между нулём и единицей. Во времени нужны метрики вроде вероятности компрометации за заданный интервал.
  \item Аналитические модели, поведенческие симуляции, сетевая эмуляция или симуляция и измерения публичной сети отвечают на разные вопросы и несут разные риски.
  \item В симуляции по эталонному консенсусу постоянные guard существенно уменьшают долю клиентов, компрометируемых со временем, концентрируя экспозицию у меньшего числа пользователей.
\end{itemize}

% ---------------------------------------------------------------------------
\section*{Упражнения}

\begin{esercizio}[1]\label{es:29-entropia}
Атакующий наблюдает систему из четырёх пользователей и назначает вероятности $1/2$, $1/4$, $1/8$ и $1/8$. Вычислите энтропию, степень анонимности и эквивалентное число равновероятных пользователей.
\end{esercizio}

\begin{esercizio}[1]\label{es:29-metodo}
Выберите и обоснуйте наиболее подходящий метод для каждого вопроса: (a) как изменилась бы задержка, если бы цепи использовали четыре перехода? (b) Сколько пользователей подключается к Tor из конкретной страны? (c) Как новый алгоритм выбора сторожевого узла изменил бы время до первого компрометационного события?
\end{esercizio}

\begin{esercizio}[2]\label{es:29-durata}
Запустите \texttt{simulate} из Лабораторной работы~19 с \texttt{guard\_lifetime}, равным 30 и 180 дням вместо значения по умолчанию 60, и сравните долю скомпрометированных клиентов через 180 дней. Объясните результаты.
\end{esercizio}

\begin{esercizio}[2]\label{es:29-asimmetria}
Повторите Лабораторную работу~19 с атакующим, контролирующим 10\% пропускной способности выходных узлов, но не сторожевых, а затем с атакующим, контролирующим 10\% пропускной способности сторожевых узлов, но не выходных. Чего способен добиться каждый и почему?
\end{esercizio}

\begin{esercizio}[3]\label{es:29-torps}
Расширьте симулятор тремя функциями, которые текущая модель игнорирует: повторное использование клиентами цепей в течение десяти минут вместо постоянного создания новых, ограничение /16 между guard и exit и атакующий, контролирующий целую автономную систему крупной сети. Какие дополнения сильнее всего меняют результат?
\end{esercizio}

% ---------------------------------------------------------------------------
\section*{Дополнительная литература}
\begin{itemize}
  \item C.~Díaz, S.~Seys, J.~Claessens, B.~Preneel, \titolo{Towards Measuring Anonymity}~\cite{diaz2002measuring}.
  \item R.~Jansen, A.~Johnson, \titolo{Safely Measuring Tor}~\cite{jansen2016privcount}.
  \item Исследовательский портал Tor Project: \url{https://research.torproject.org/}.
  \item \emph{Anonymity Bibliography}: \url{https://www.freehaven.net/anonbib/}.
\end{itemize}
